% !TeX program = pdflatex
%-------------------------
% Lettre de motivation
% Author : Javier Tarazona
% Offre : BNP Paribas - Stage - AI Research Intern H/F (ITG Research, ref S_ITGIDF27_0088)
%------------------------

\documentclass[a4paper,11pt]{article}
\usepackage[T1]{fontenc}
\usepackage[utf8]{inputenc}
\usepackage{lmodern}
\usepackage[french]{babel}
\usepackage[a4paper,top=1.4cm,bottom=1.4cm,left=2cm,right=2cm]{geometry}
\usepackage[hidelinks]{hyperref}
\usepackage{microtype}

\hypersetup{
  pdftitle={Javier Andres TARAZONA JIMENEZ - Lettre de motivation - AI Research Intern - BNP Paribas},
  pdfauthor={Javier Andres TARAZONA JIMENEZ}
}

\pagestyle{empty}
\setlength{\parindent}{0pt}
\setlength{\parskip}{5pt}

\begin{document}

%----------EXPEDITEUR-----------------
{\bfseries Javier Andres TARAZONA JIMENEZ}\\
Palaiseau, France\\
+33 7 46 36 97 75\\
\href{mailto:javitar06@gmail.com}{javitar06@gmail.com}

%----------DESTINATAIRE-----------------
\hfill\begin{minipage}[t]{0.42\textwidth}
{\bfseries BNP Paribas}\\
ITG Research\\
À l'attention de l'équipe de recrutement\\
Montreuil, France
\end{minipage}

\vspace{4pt}
\hfill Palaiseau, le 7 octobre 2026

\vspace{2pt}
\textbf{Objet : Candidature au stage « AI Research Intern H/F »}

Madame, Monsieur,

Étudiant en double diplôme entre l'Universidad Nacional de Colombia et l'ENSTA Paris, école de l'Institut Polytechnique de Paris, je suis venu à l'informatique par curiosité pour des systèmes dont la simplicité apparente cache une grande complexité -- Internet, infrastructures financières, caméras -- et par l'envie de transformer des idées abstraites en systèmes que je construis moi-même. L'intelligence artificielle a ensuite élargi cette approche : plutôt que de programmer chaque comportement, concevoir des systèmes qui extraient l'information des données, modélisent le réel et éclairent la décision. Au sein du parcours IA de l'ENSTA, cet intérêt s'est concrétisé selon trois axes : l'analyse de données, au laboratoire U2IS et pour l'aide à la décision du bateau autonome « Tameo » ; la vision et le langage, avec le pipeline de détection et de décision de Tameo, l'analyse de marquages routiers chez Engin A.I et une application de NLP pour enfants dyslexiques chez APEDYS91 ; enfin le développement logiciel, car l'IA ne crée de valeur qu'une fois déployée dans un service fonctionnel, comme chez Engin A.I.

ITG Research m'attire d'abord par le terrain qu'offre BNP Paribas. Banque de financement et d'investissement, banque commerciale et assurance : ces trois métiers produisent des données de marchés, d'actifs et de risques qui appellent analyse et décision. Plus de 800 cas d'usage d'IA en production, de nombreux data scientists, des liens académiques -- notamment avec Télécom Paris --, des chercheurs publiant à NeurIPS ou ICSE et le partenariat avec Mistral AI montrent une recherche reliée à des usages réels. Les thématiques du stage -- foundation models, LLM, IA de confiance, IA pour la finance -- rejoignent la direction que je veux donner à mon parcours.

Mon stage au laboratoire U2IS recoupe plusieurs missions du poste : j'y ai pré-entraîné des Vision Transformers en PyTorch par apprentissage auto-supervisé, conçu un protocole en curriculum learning et construit un pipeline Python reproductible de génération de données multi-vues. J'en retiens une compréhension concrète de l'apprentissage de représentations, au cœur des foundation models. Chez APEDYS91, j'ai combiné règles et prompting de LLM locaux (Mistral, Gemma), avec des garde-fous contre les hallucinations et un choix de modèle contraint par la mémoire disponible : un premier contact avec les enjeux de fiabilité et de coût de calcul. Sur Tameo, dont je suis désormais chef de projet, le fine-tuning et la comparaison de détecteurs (mAP et F1 portés d'environ 80\,\% à 90\,\%) m'ont appris à bâtir un benchmark face à des références. Le RAG et la finance quantitative relèvent encore de mes centres d'intérêt plus que de mon expérience ; je les aborderai avec mes bases en probabilités, statistiques et modèles stochastiques. Mon parcours entre la Colombie, où le Groupe est présent, le Canada et Paris m'a enfin habitué aux environnements internationaux.

En retour, ce stage prolongerait mon expérience de laboratoire vers l'industrie : passer de l'état de l'art et des hypothèses à des prototypes soumis aux contraintes de robustesse et de passage à l'échelle d'un grand groupe, dans la continuité de mes expériences produit chez Engin A.I et APEDYS91, avec l'encadrement d'équipes expérimentées. Les travaux du Groupe, du FinAI-Lab sur l'apprentissage continu sur de grands réseaux de données évolutifs aux foundation models pour séries temporelles financières et à la modélisation générative par ponts de Schrödinger, que je commence à découvrir, illustrent les problèmes auxquels je souhaite contribuer. La perspective d'une thèse CIFRE, si les résultats s'y prêtent, m'intéresse particulièrement. Je serais heureux d'échanger avec vous sur cette candidature lors d'un entretien.

Je vous prie d'agréer, Madame, Monsieur, l'expression de mes salutations distinguées.

\hfill Javier Andres TARAZONA JIMENEZ

\end{document}
